\chapter{Introducción}
\label{cap:introduccion}

\section{Contexto}
\label{contexto}

La soledad forma parte de la vida y, cuando se elige, puede resultar incluso reparadora. La dificultad aparece cuando no se elige, sino que viene impuesta por circunstancias sociales, demográficas, familiares o laborales. Esa soledad no deseada (en adelante, SND) afecta de manera particular a la tercera edad, donde suele venir precedida por el fallecimiento de la pareja, por la reducción del núcleo familiar o por el distanciamiento de los círculos sociales de referencia. \textcite{RuedaEstrada2018} la describe como un problema de preocupación creciente en España, que ha dejado de interpelar solo a la ciudadanía para hacerlo también a geriatras, psicólogos, trabajadores sociales y responsables de la gestión pública.

El problema tiene ya reconocimiento institucional. En 2024 la \textit{Revista Española de Salud Pública}, editada por el Ministerio de Sanidad, describía la SND como \enquote{un problema silente de nuevo orden} \citep{JimenezRodriguez2024}, y el \textit{Barómetro de la soledad no deseada en España}, del Observatorio Estatal impulsado por Fundación ONCE y Fundación AXA, la sitúa en el 20\% de la población adulta \citep{SoledadES}. Entre los mayores de 65 años son cerca de tres millones de personas los que la padecen, con un coste para el Estado estimado en 14.141 millones de euros anuales \citep{Senara}. Lo característico de la soledad en la vejez, y lo que la distingue de la que se sufre a otras edades, es que suele estar ligada a pérdidas difícilmente reversibles y tiende por ello a cronificarse \citep{FundacionONCE}. El capítulo~\ref{cap:marco-teorico} desarrolla la dimensión del problema y la evidencia sobre sus consecuencias.

Sobre ese problema se superpone una barrera tecnológica. La digitalización de la banca, la administración y la sanidad ha convertido la competencia digital en requisito para la participación social justo cuando se cierran los canales presenciales: el 57\% de los municipios españoles carece ya de oficina bancaria \citep{RevistaByte}, y más de un tercio de la población queda desconectada del entorno digital por carecer de dispositivos, de una conexión adecuada o de las habilidades para usarlos, sobre todo a partir de los 65~años \citep{espanol}. Las herramientas que podrían mitigar la soledad quedan así fuera del alcance de quienes más las necesitan.

Y esa barrera no es solo de capacitación, sino también de diseño del propio software. Las aplicaciones de uso corriente dan por supuesta una alfabetización digital que este colectivo no tiene: áreas táctiles pequeñas, tipografías reducidas y gestos sin afordancia visible, como el deslizamiento o la pulsación mantenida, que son precisamente los que más fallan en usuarios mayores \citep{Kane2019}. Las que sí se dirigen a ellos se reparten entre dos extremos igual de problemáticos: por un lado el hardware dedicado, de coste elevado e implantación testimonial en España; por otro los compañeros conversacionales generalistas, que no se diseñaron para este perfil sino que se optimizan para maximizar el tiempo de uso \citep{Zhou2020}, con patrones documentados de dependencia emocional y expedientes regulatorios abiertos en Europa.

Esa intersección entre una necesidad social contrastada y una oferta de software que no la atiende es la que justifica este Trabajo Fin de Grado.

\section{Motivación}

La motivación de este proyecto es doble. En el plano técnico se trata de completar el desarrollo de una aplicación con decisiones justificadas, mientras se ponen en práctica las competencias adquiridas a lo largo del grado. Merece la pena nombrar cuáles, porque explican buena parte de lo que documenta esta memoria.
\begin{itemize}
    
    \item \textit{Desarrollo de Software}: de esta asignatura destaco el descubrimiento de Flutter como \emph{framework}, con el que se hicieron varias aplicaciones pequeñas, y los patrones de diseño. De Flutter se aprovechan dos rasgos: la interfaz se compone de \emph{widgets} reutilizables, y sus paquetes dan acceso al reconocimiento y a la síntesis de voz sin programar contra las API nativas de cada plataforma. De los patrones de diseño, el proyecto aplica tres. El primero es la inyección de dependencias: cada clase recibe desde fuera los servicios que usa, en lugar de crearlos ella misma, de modo que en las pruebas pueden sustituirse por versiones simuladas (apartado~\ref{sec:estado-di}). Los otros dos son el repositorio, que aísla el acceso a datos, y el observador, con el que la interfaz se actualiza sola cuando cambia el perfil del usuario.

    \item \textit{Metodologías de Desarrollo Ágil} y \textit{Dirección y Gestión de Proyectos}: cómo elegir metodología según las condiciones del proyecto y cómo pasar de las necesidades de un cliente a historias de usuario, de estas a requisitos y de ahí al código. Sus proyectos finales, junto con el de \textit{Desarrollo de Aplicaciones para Internet}, consistieron en aplicaciones web y móviles hechas con esas metodologías, en equipo y en solitario. El capítulo~\ref{cap:metodologia} las adapta a un proyecto de una sola persona, conservando solo las prácticas que tienen sentido sin equipo. 

    \item \textit{Programación y Diseño Orientado a Objetos} y \textit{Fundamentos de Ingeniería del Software}: el modelado con diagramas de casos de uso, de clases y de secuencia antes de escribir código, porque un error de diseño se corrige antes en un diagrama que en la implementación. Los capítulos~\ref{cap:analisis} y~\ref{cap:diseno-desarrollo} recogen esos modelos, junto con el diagrama de arquitectura global del sistema.

    \item \textit{Fundamentos del Software} y \textit{Sistemas Operativos}: algo menos vistoso pero de uso diario, la soltura con la línea de comandos. Todo el ciclo de trabajo de este proyecto vive ahí, desde el control de versiones y las herramientas de construcción hasta el despliegue de las funciones en la nube, la ejecución de las suites de pruebas y la lectura de los registros del dispositivo cuando algo falla.

\end{itemize}

En el plano social, el producto busca ocupar el hueco descrito en el apartado anterior, el que queda entre un hardware dedicado que casi nadie puede pagar y unos compañeros conversacionales que no se diseñaron para personas mayores. No pretende, en cambio, ser una intervención terapéutica contra la soledad, pues no se ha realizado ningún estudio clínico con usuarios reales que permita afirmarlo.

Lo que se propone es una tercera vía: una aplicación para un teléfono móvil corriente, con un coste de despliegue mínimo, diseñada en torno a la accesibilidad y en español, que complementa el vínculo humano en lugar de sustituirlo. Esa premisa se materializa en el modelo dual, que vincula a la persona mayor con un familiar y es a la vez el rasgo diferencial de la propuesta y su respuesta a los riesgos de las alternativas comerciales, analizadas en el capítulo~\ref{cap:marco-teorico}.

\section{Objetivos}

La propuesta, Quiet Companion, parte de tres principios de diseño: una interfaz accesible, la voz como forma principal de interacción y un cuidado compartido con un familiar. Este último se concreta en la \textit{revelación progresiva}: la aplicación se adapta al grado de autonomía de cada usuario. Los tres se desarrollan en el capítulo~\ref{cap:diseno-desarrollo}.

El objetivo general del trabajo es \textbf{desarrollar un prototipo funcional de aplicación móvil que acompañe por voz a personas mayores en situación de soledad no deseada, accesible para ese perfil de usuario}. Se desglosa en los objetivos específicos de la tabla~\ref{tab:objetivos}: los cinco primeros se refieren al producto y el último, al proceso de ingeniería.

{\setlength{\intextsep}{6pt}% menos aire alrededor: así la tabla cabe en la página de su apartado
\begin{table}[H]
    \centering
    \caption{Objetivos específicos del proyecto.}
    \label{tab:objetivos}
    \small
    \begin{tabularx}{\textwidth}{@{}c >{\RaggedRight\arraybackslash}X@{}}
        \toprule
        \textbf{ID} & \textbf{Descripción} \\
        \midrule
        OE-1 & Implementar un sistema de gestión de perfiles con revelación progresiva que permita al usuario operar en modo independiente o vinculado a un familiar, así como el mecanismo de emparejamiento entre ambos. \\
        \addlinespace[2pt]
        OE-2 & Desarrollar una interfaz conversacional por voz que integre el ciclo completo de reconocimiento de voz, generación de respuesta mediante un modelo de lenguaje y síntesis de voz, comunicando visualmente sus estados, y ofrecer la escritura como vía alternativa para quien no pueda o no quiera hablar. \\
        \addlinespace[2pt]
        OE-3 & Cumplir los criterios de accesibilidad de la norma WCAG 2.2 en la interfaz dirigida a la persona mayor, mediante tipografías de gran tamaño, áreas táctiles amplias e interacción simplificada. \\
        \addlinespace[2pt]
        OE-4 & Implementar un sistema de configuración adaptativa que permita ajustar los parámetros de accesibilidad y que propague en tiempo real los cambios realizados por el familiar al dispositivo del mayor. \\
        \addlinespace[2pt]
        OE-5 & Incorporar el seguimiento del bienestar del mayor, mediante una memoria conversacional y el análisis del estado de ánimo de cada conversación, supeditado al consentimiento informado del propio mayor conforme al RGPD y sin almacenar en ningún momento la transcripción de lo hablado. \\
        \addlinespace[2pt]
        OE-6 & Diseñar e implementar la solución sobre una arquitectura escalable y mantenible, aplicando los principios de Clean Architecture y Feature-Sliced Design, y gestionar el ciclo de vida del proyecto mediante una metodología ágil basada en tableros Kanban. \\
        \bottomrule
    \end{tabularx}

    \vspace{2pt}
    {\footnotesize Fuente: Elaboración propia.}
\end{table}}

\section{Estructura de la memoria}

Además de esta introducción, la memoria consta de cinco capítulos. El capítulo~\ref{cap:marco-teorico} revisa el marco teórico de la soledad no deseada y el estado del arte de las soluciones tecnológicas. El capítulo~\ref{cap:metodologia} justifica la metodología y recoge la planificación y el estudio de costes. El capítulo~\ref{cap:analisis} especifica actores, requisitos, casos de uso y modelo de dominio. El capítulo~\ref{cap:diseno-desarrollo} documenta el diseño y el desarrollo. El capítulo~\ref{cap:desarrollo-futuro} cierra con las conclusiones y las líneas de trabajo futuro, y el apéndice~\ref{ap:manual} recoge el manual de usuario.

La aplicación y su código fuente están en el repositorio:
\newline
\url{https://github.com/maario106/quiet_companion}
